\documentclass[letterpaper,12pt]{article}
\usepackage[utf8]{inputenc}
\usepackage{charter}
\usepackage{geometry}
\usepackage{amsmath}
\usepackage{float}
\usepackage{graphicx}
\usepackage{subcaption}
\usepackage{amssymb}
\usepackage{adjustbox}
\usepackage{wrapfig} 
\usepackage{xcolor}
\usepackage{fancyhdr}
\usepackage{tabularx}
\usepackage{fancyhdr}
\usepackage{comment}

\geometry{top=2cm, bottom=2cm, left=2cm, right= 2cm} %%margen
\graphicspath{{images/}}
\parindent=0pt

\begin{document}
%%para contador y eso
\thispagestyle{empty}
\newpage
\setcounter{page}{1}
\pagestyle{headings}
%%%%%%%%%%%%%%%%%%%%%%%%%%%%%%%%%%%%%%%%%%%%%%%%%%%%%%%%%%%%%%%%%%%%%%%%%%%%
\begin{sloppypar} 
    %%%portadita
    \begin{titlepage}
        \fancyhead{R}{
        \includegraphics[width=0.1\linewidth]{logoUV.png}
        }
        \hspace{2.5cm}
        {\bfseries\LARGE Universidad Veracruzana \par}
        \hspace{2cm}
        {\scshape\Large Facultad de Ingeniería Eléctrica y Electrónica \par}
        \begin{center}
            \vspace{7cm}
            {\itshape\huge Actividades de Métodos Numéricos \par}
            {\large Lara Xocuis Martha Denisse\par}
            {\large S22002213 \par}
            \vfill
            {\Large 27 de Agosto, 2024 \par}
        \end{center}
    \end{titlepage}
    
\section*{Tarea 1}
\begin{center}
    \begin{tabular}[H]{|c|c|}\hline 
        $v(t)$ & $t$ \\ \hline
        0 & 0 \\ \hline
        16.43 & 2 \\ \hline
        27.85 & 4 \\ \hline
        35.77 & 6 \\ \hline
        41.27 & 8 \\ \hline
        45.09 & 10 \\ \hline
        47.74 & 12  \\ \hline
        49.58 & 14 \\ \hline
        50.85 & 16 \\ \hline
    \end{tabular} 
\end{center}

$\displaystyle v(t)  = \frac{gm}{c} (1 - e^{-(c/m)t})$

\begin{center}
    \begin{tabular}[H]{|c|c|}\hline 
        $v_{i+1}$ & $t_{i+1}$ \\ \hline
        0 & 0 \\ \hline
        19.62 & 2 \\ \hline
        32.07 & 4 \\ \hline
        39.98 & 6 \\ \hline
        45.00 & 8 \\ \hline
        48.19 & 10 \\ \hline
        50.17 & 12  \\ \hline
        51.45 & 14 \\ \hline
        52.27 & 16 \\ \hline
    \end{tabular} 
\end{center}


\end{sloppypar}
\end{document}